% =============================================================================
% Informe Técnico: Detección de Vishing mediante Biometría Comportamental
% Formato: IEEE Transactions (journal, doble columna)
% Instrucciones para figuras: busca los comentarios %[FIGURA] para insertar
%   las imágenes exportadas de los notebooks. Ruta sugerida indicada en cada caso.
% =============================================================================
\documentclass[journal]{IEEEtran}

% ── Idioma y codificación ─────────────────────────────────────────────────────
\usepackage[spanish,es-tabla]{babel}
\usepackage[utf8]{inputenc}
\usepackage[T1]{fontenc}

% ── Matemáticas ───────────────────────────────────────────────────────────────
\usepackage{amsmath}
\usepackage{amssymb}

% ── Gráficos y tablas ─────────────────────────────────────────────────────────
\usepackage{graphicx}
\usepackage{booktabs}
\usepackage{array}
\usepackage{multirow}
\usepackage{tabularx}
\usepackage{xcolor}
\usepackage{float}

% ── TikZ (diagrama de pipeline) ───────────────────────────────────────────────
\usepackage{tikz}
\usetikzlibrary{shapes.geometric, arrows.meta, positioning, fit, backgrounds, calc}

% ── Referencias y URLs ────────────────────────────────────────────────────────
\usepackage{cite}
\usepackage{url}
\usepackage{hyperref}
\hypersetup{
    colorlinks=true,
    linkcolor=blue!60!black,
    citecolor=blue!60!black,
    urlcolor=blue!60!black,
}

% ── Algoritmos ────────────────────────────────────────────────────────────────
\usepackage{algorithm}
\usepackage{algorithmic}

% ── Comandos auxiliares ───────────────────────────────────────────────────────
\newcommand{\wrapcls}{\texttt{VishingModelWrapper}}
\newcommand{\augcls}{\texttt{CorrelatedAugmenter}}
\newcommand{\prauc}{\textsc{PR-AUC}}
\newcommand{\rocauc}{\textsc{ROC-AUC}}

% =============================================================================
\begin{document}

\title{Detección de Sesiones de Vishing en Banca Digital\\
Mediante Biometría Comportamental y Aprendizaje Automático:\\
Una Prueba de Concepto con Datos Sintéticos}

\author{%
  \IEEEauthorblockN{[Nombre Autor 1], [Nombre Autor 2]}\\
  \IEEEauthorblockA{[Unidad / Área]\\
    [Entidad Bancaria]\\
    Bogotá, Colombia\\
    \texttt{[correo@institucion.com]}}
}

\maketitle

% =============================================================================
\begin{abstract}
El vishing (\textit{voice phishing}) es una modalidad de fraude en la que un
atacante guía telefónicamente a un usuario bancario para que ejecute
transacciones en su nombre. A diferencia del phishing tradicional, el usuario
opera voluntariamente dentro de la aplicación legítima, lo que anula las
defensas habituales basadas en canal o dispositivo. Este trabajo presenta un
pipeline completo de machine learning para detectar sesiones de vishing en
tiempo real a partir de señales de biometría comportamental inspiradas en el
marco tecnológico de BioCatch. Ante la ausencia de datos etiquetados
disponibles, se construyó un dataset sintético de \num{50000} sesiones y se
amplió a \num{1000000} mediante \augcls{}, un generador que preserva la
estructura de correlación multivariada mediante descomposición de Cholesky y
transformación cuantil inversa. Se evaluaron cuatro algoritmos sobre doce
configuraciones de balanceo de clases y posteriormente se realizó una
experimentación dirigida con siete variantes hiperparamétricas de XGBoost
sobre veintiseis datasets, totalizando 182~modelos. El mejor resultado en
holdout alcanzó \prauc{}$=0.9484$, \textit{Recall}$=0.878$ y
$F_1=0.902$. La simulación de inferencia en tiempo real sobre 100\,000
sesiones sintéticas adicionales produjo un \textit{Recall}$=0.755$,
$F_1=0.763$ y una latencia de inferencia media de $1.60$~ms, con
percentil~99 en $2.24$~ms y un rendimiento de $539$~observaciones/s,
valores consistentes con despliegue en producción. Todos los modelos se
serializaron como \wrapcls{} y se almacenan en Amazon S3.
\end{abstract}

\begin{IEEEkeywords}
vishing, biometría comportamental, detección de fraude, aprendizaje automático,
XGBoost, datos sintéticos, desbalance de clases, banca digital, BioCatch,
inferencia en tiempo real.
\end{IEEEkeywords}

% =============================================================================
\section{Introducción}
\label{sec:intro}

\IEEEPARstart{L}{a} adopción masiva de la banca digital ha desplazado el
fraude hacia canales digitales, donde el vector humano sigue siendo el eslabón
más vulnerable. El vishing combina ingeniería social por voz con la interfaz
legítima de una aplicación bancaria: el atacante llama a la víctima, se
hace pasar por un empleado del banco o una autoridad, y la instruye paso a
paso para transferir fondos a cuentas controladas por el delincuente
\cite{mitnick2002art, jagatic2007social}.

La dificultad central del problema radica en que el usuario \emph{sí} es
quien opera la aplicación, desde su dispositivo habitual y con sus
credenciales válidas. Los sistemas de autenticación multifactor, detección
de dispositivos nuevos o validación de IP fallan porque todas las señales
de identidad resultan consistentes. Lo que difiere es \emph{cómo} el usuario
interactúa: escribe más lento porque dicta lo que le indica el atacante,
hace pausas largas mientras escucha instrucciones, comete más errores de
entrada al transcribir datos que no memoriza, y aumenta los montos de
transacción de forma inusual~\cite{threatfabric2023}.

La biometría comportamental captura precisamente estas diferencias
\cite{shen2007application, sundararajan2020behavioral}. Empresas como
BioCatch han construido plataformas comerciales que analizan en tiempo real
cientos de señales de comportamiento durante una sesión bancaria para
detectar patrones anómalos asociados a fraude, ingeniería social y acceso
remoto no autorizado~\cite{biocatch2022}.

Este trabajo documenta una prueba de concepto interna que busca responder:
\emph{¿es posible entrenar un clasificador de sesiones de vishing a partir de
biometría comportamental sintética, con métricas y latencia aptas para
producción?}

\subsection*{Contribuciones principales}

\begin{enumerate}
    \item \textbf{Dataset sintético de biometría comportamental}: \num{50000}
    sesiones (47\,500 legítimas, 2\,500 de vishing) con 61~variables
    diseñadas a partir de supuestos públicos de BioCatch.

    \item \textbf{\augcls{}}: aumentador correlacionado que expande el
    dataset a $1\times10^6$ sesiones preservando la estructura de correlación
    multivariada mediante descomposición de Cholesky y emparejamiento
    cuantil, superando las limitaciones de SMOTE en datos tabulares
    comportamentales.

    \item \textbf{Pipeline comparativo de cuatro fases}: evaluación sistemática
    de cuatro algoritmos, doce estrategias de balanceo y siete variantes de
    XGBoost sobre dos tipos de datos (50K y 1M), totalizando 230~modelos.

    \item \textbf{\wrapcls{}}: patrón de empaquetado para producción que
    encapsula modelo, lista de features y umbral óptimo en un único artefacto
    serializable.

    \item \textbf{Simulación de inferencia en tiempo real}: validación del
    pipeline de punta a punta sobre 100\,000~sesiones 100\,\% sintéticas,
    demostrando viabilidad operacional ($\bar{x}_{\text{lat}}=1.60$~ms).
\end{enumerate}

El resto del documento se organiza así: la Sección~\ref{sec:arte} revisa el
estado del arte; la Sección~\ref{sec:dataset} describe el dataset; la
Sección~\ref{sec:metodologia} detalla la metodología; la
Sección~\ref{sec:wrapper} describe el wrapper de producción; la
Sección~\ref{sec:resultados} presenta los resultados; la
Sección~\ref{sec:simulacion} analiza la simulación en tiempo real; la
Sección~\ref{sec:discusion} discute hallazgos y limitaciones; y la
Sección~\ref{sec:conclusiones} cierra el documento.

% =============================================================================
\section{Estado del Arte}
\label{sec:arte}

\subsection{Biometría comportamental en banca digital}

La biometría comportamental analiza patrones dinámicos de interacción
---dinámica de teclado, presión táctil, velocidad de scroll, inclinación del
dispositivo--- que son difíciles de falsificar porque emergen de hábitos
motores inconscientes~\cite{fridman2017active, sundararajan2020behavioral}.
A diferencia de la biometría fisiológica (huella dactilar, rostro), no
requiere hardware especializado y opera de forma pasiva durante la sesión.

En el contexto bancario, estas señales se han aplicado exitosamente para
detectar bots, ataques de tipo \textit{account takeover} (ATO) y acceso
remoto no autorizado~\cite{biocatch2022}. La detección de vishing es un
caso más complejo: el atacante no controla el dispositivo directamente, sino
que actúa a través del comportamiento inducido en la víctima.

\subsection{Detección de fraude con aprendizaje automático}

Los enfoques supervisados dominan la detección de fraude bancario. Las
revisiones sistemáticas de Phua et al.~\cite{phua2010comprehensive} y
Bhattacharyya et al.~\cite{bhattacharyya2011data} muestran que los
ensambles de árboles de decisión ---Random Forest y gradient
boosting--- ofrecen el mejor balance entre rendimiento y velocidad de
inferencia para aplicaciones en tiempo real.

XGBoost~\cite{chen2016xgboost} se ha convertido en el estándar de facto para
datos tabulares en fraude por su manejo implícito de valores atípicos,
regularización incorporada y eficiencia en GPU. En competencias de referencia
como las de Kaggle sobre detección de fraude en tarjetas de crédito, las
soluciones ganadoras invariablemente incluyen XGBoost o LightGBM como
componente central.

El desbalance de clases es un desafío estructural: las tasas de fraude reales
oscilan entre el 0.1\,\% y el 5\,\%, lo que sesga los clasificadores hacia
la clase mayoritaria~\cite{he2009learning, krawczyk2016learning}. Las técnicas
de sobre-muestreo como SMOTE~\cite{chawla2002smote} y sus variantes
(Borderline-SMOTE~\cite{han2005borderline}) son el enfoque más utilizado,
aunque generan muestras sintéticas por interpolación lineal sin preservar
correlaciones de orden superior.

\subsection{Ataques de vishing e ingeniería social}

El vishing es una variante de la ingeniería social~\cite{tetri2013dissecting}
en la que el canal de ataque es la voz. A diferencia del phishing por correo
electrónico, el atacante adapta el guion en tiempo real según las respuestas
de la víctima, lo que aumenta la tasa de éxito~\cite{jagatic2007social}.

Los ataques de vishing bancario suelen incluir: urgencia artificial (amenaza
de bloqueo de cuenta), suplantación de identidad (número de teléfono
falso), y guion paso a paso para completar transacciones~\cite{apwg2023}.
La detección basada únicamente en el canal de voz (análisis de llamadas) es
insuficiente porque requiere acceso al audio y opera fuera del perím-etro de
la aplicación bancaria.

\subsection{Datos sintéticos para entrenamiento de modelos de fraude}

La escasez de datos etiquetados es un obstáculo reconocido en la
investigación de fraude: los datos reales son escasos, confidenciales y
sujetos a restricciones regulatorias~\cite{patki2016synthetic}. La generación
de datos sintéticos que preserven las distribuciones estadísticas y las
correlaciones del mundo real es un campo activo, con enfoques que van desde
modelos generativos adversariales (GANs)~\cite{nikolentzos2021synthetic}
hasta métodos más interpretables basados en estadística clásica.

El enfoque adoptado en este trabajo ---preservar la estructura de correlación
multivariada mediante descomposición de Cholesky y transformación cuantil--- es
similar al propuesto en el marco Synthetic Data Vault~\cite{patki2016synthetic},
pero adaptado a los requisitos específicos de variables comportamentales con
restricciones de dominio (no negatividad, rangos acotados, variables
binarias dependientes de otras variables de sesión).

% =============================================================================
\section{Dataset Sintético}
\label{sec:dataset}

\subsection{Proceso de construcción}

El dataset fue construido con el apoyo de IA generativa (Claude, Anthropic) a
partir de información pública de BioCatch sobre vectores de comportamiento
asociados a vishing~\cite{biocatch2022}. BioCatch no divulga las variables
exactas de su plataforma, pero describe en alto nivel los tipos de señal que
captura y los patrones comportamentales asociados a ingeniería social. A partir
de estos supuestos se diseñaron 61~variables que cubren ocho grupos
funcionales.

La generación controlada con parámetros estadísticos distintos por clase
(legítima vs. vishing) permite asegurar que los patrones simulados sean
coherentes con las hipótesis comportamentales del problema: sesiones de vishing
exhiben tipeo más lento, mayor hesitación, más tiempo muerto y montos de
transacción significativamente más altos.

\subsection{Composición y período simulado}

El dataset original comprende \num{50000} sesiones simuladas para el período
junio 2024 -- mayo 2025, con una tasa de vishing del 5\,\% (Tabla~\ref{tab:dataset}).

\begin{table}[h]
\caption{Composición del dataset de biometría comportamental sintética.}
\label{tab:dataset}
\centering
\begin{tabular}{lrr}
\toprule
\textbf{Conjunto} & \textbf{Sesiones} & \textbf{\% Vishing} \\
\midrule
Dataset original (50K)          & \num{50000}   & 5.00\,\% \\
\quad Legítimas                 & \num{47500}   & ---      \\
\quad Vishing                   & \num{2500}    & ---      \\
\midrule
Dataset aumentado (1M)          & \num{1000000} & 1.50\,\% \\
\quad Legítimas (orig.\ + sint.)& \num{985000}  & ---      \\
\quad Vishing (orig.\ + sint.)  & \num{15000}   & ---      \\
\midrule
Dataset inferencia (100K sint.) & \num{100000}  & 1.50\,\% \\
\bottomrule
\end{tabular}
\end{table}

\subsection{Variables comportamentales}

Se definieron 61~variables organizadas en diez grupos; se retienen 44~para
el modelado tras eliminar identificadores y variables de BioCatch que
generarían \textit{data leakage} (Tabla~\ref{tab:variables}).

\begin{table}[h]
\caption{Grupos de variables comportamentales del dataset.}
\label{tab:variables}
\centering
\begin{tabularx}{\columnwidth}{lXr}
\toprule
\textbf{Grupo} & \textbf{Señal capturada} & \textbf{Vars.} \\
\midrule
Dinámica de teclado   & Velocidad, latencia entre teclas, variabilidad, ratio de escritura segmentada & 5 \\
Dinámica táctil       & Presión, tamaño de toque, velocidad y varianza de swipes & 5 \\
Movimiento dispositivo& Ángulo, giroscopio, acelerómetro, eventos de movimiento & 5 \\
Hesitación            & Cantidad y duración de pausas (usuario espera instrucciones) & 3 \\
Tiempo muerto         & Períodos de inactividad (usuario escucha al atacante) & 3 \\
Navegación en app     & Pantallas únicas visitadas, retrocesos, tiempo de transición & 3 \\
Errores y correcciones& Errores de entrada, correcciones de monto/beneficiario, copiar/pegar & 4 \\
Contexto de sesión    & Duración, hora, llamada activa, acceso remoto, apps sospechosas & 5 \\
Datos de transacción  & Monto (COP), beneficiario nuevo, tiempo hasta transacción & 4 \\
Features derivadas    & Ratio de tiempo muerto, hesitación compuesta, interacciones/s & 7 \\
\midrule
\textbf{Total para modelado} & \multicolumn{1}{r}{} & \textbf{44} \\
\bottomrule
\end{tabularx}
\end{table}

La variable \texttt{phone\_call\_active} merece mención especial: indica si
hay una llamada telefónica activa durante la sesión bancaria, y es la señal
más discriminativa del dataset. La aplicación bancaria de referencia sí tiene
acceso a este estado a través de las APIs del sistema operativo.

% =============================================================================
\section{Metodología}
\label{sec:metodologia}

\subsection{Visión general del pipeline}

El pipeline se organiza en cuatro fases iterativas que progresivamente aumentan
la escala de datos y la profundidad del espacio de hiperparámetros explorado
(Fig.~\ref{fig:pipeline}). Las fases~1 y el análisis EDA de la fase~2 se
ejecutan en entorno local; las fases de augmentación y entrenamiento a escala
se ejecutan en AWS SageMaker con GPU (\texttt{ml.g4dn.xlarge} o superior).

% ── DIAGRAMA DE PIPELINE (TikZ — no requiere imagen externa) ─────────────────
\begin{figure*}[t]
\centering
\begin{tikzpicture}[
    pbox/.style={rectangle, rounded corners=5pt, draw=blue!55, fill=blue!8,
                 text width=3.1cm, minimum height=3.8cm, align=center,
                 font=\small, inner sep=5pt},
    rbox/.style={rectangle, rounded corners=3pt, draw=orange!70, fill=orange!12,
                 text width=2.9cm, minimum height=0.85cm, align=center,
                 font=\footnotesize\bfseries, inner sep=3pt},
    dbox/.style={rectangle, rounded corners=3pt, draw=green!60!black, fill=green!8,
                 text width=2.5cm, minimum height=0.7cm, align=center,
                 font=\scriptsize, inner sep=2pt},
    arr/.style={thick, -Latex, gray!70},
    darr/.style={thick, dashed, -Latex, gray!55},
]

%% ── Fase 1 ───────────────────────────────────────────────────────────────────
\node[pbox] (p1) at (0,0) {
    \textbf{FASE 1}\\[3pt]
    EDA \& Análisis\\
    Balanceo (12~var.)\\
    4 Algoritmos\\
    Holdout 10K\\[5pt]
    \textit{\footnotesize Local — 50K}
};
\node[dbox, above=0.35cm of p1] (d1) {50K sesiones\\5\,\% vishing};
\node[rbox, below=0.35cm of p1] (r1) {PR-AUC $\approx$ 0.45\\XGBoost mejor};

%% ── Fase 2 ───────────────────────────────────────────────────────────────────
\node[pbox, right=0.9cm of p1] (p2) at (4.3,0) {
    \textbf{FASE 2}\\[3pt]
    \augcls\\
    EDA validación\\
    Balanceo (12~var.)\\
    4 Algoritmos\\
    Holdout 200K\\[5pt]
    \textit{\footnotesize AWS — 1M}
};
\node[dbox, above=0.35cm of p2] (d2) {1M sesiones\\1.5\,\% vishing};
\node[rbox, below=0.35cm of p2] (r2) {48 modelos\\en S3};

%% ── Fase 3 ───────────────────────────────────────────────────────────────────
\node[pbox, right=0.9cm of p2] (p3) at (8.6,0) {
    \textbf{FASE 3}\\[3pt]
    7 variantes\\
    XGBoost\\
    26 datasets\\
    182 modelos\\
    \wrapcls\\[5pt]
    \textit{\footnotesize AWS — 50K+1M}
};
\node[dbox, above=0.35cm of p3] (d3) {50K + 1M\\(26 datasets)};
\node[rbox, below=0.35cm of p3] (r3) {PR-AUC = 0.9484\\F$_1$ = 0.902};

%% ── Fase 4 ───────────────────────────────────────────────────────────────────
\node[pbox, right=0.9cm of p3] (p4) at (12.9,0) {
    \textbf{FASE 4}\\[3pt]
    \augcls\\
    100K sint.\\
    Simulación\\
    tiempo real\\[5pt]
    \textit{\footnotesize AWS — 100K}
};
\node[dbox, above=0.35cm of p4] (d4) {50K como\\referencia estadística};
\node[rbox, below=0.35cm of p4] (r4) {Recall=0.755\\$\bar{x}$\,lat.=1.60\,ms};

%% ── Flechas principales ───────────────────────────────────────────────────────
\draw[arr] (p1.east) -- (p2.west);
\draw[arr] (p2.east) -- (p3.west);
\draw[arr] (p3.east) -- (p4.west);

%% ── Flechas de datos/resultados ───────────────────────────────────────────────
\foreach \n in {p1,p2,p3,p4}{
    \draw[darr] ([yshift=2pt]\n.north) -- ([yshift=-2pt]d\n |- \n.north);%
}
\foreach \n in {p1,p2,p3,p4}{
    \draw[darr] ([yshift=-2pt]\n.south) -- ([yshift=2pt]r\n |- \n.south);%
}

\end{tikzpicture}
\caption{Visión general del pipeline de cuatro fases. Cada fase amplía la
escala o el espacio de búsqueda respecto a la anterior. Los recuadros
naranjas muestran el resultado clave de cada fase.}
\label{fig:pipeline}
\end{figure*}

\subsection{Análisis exploratorio de datos (Fase 1)}

El EDA sobre el dataset de 50K sesiones estableció la viabilidad del problema
y orientó la selección de features. Se aplicaron tests de Mann-Whitney~U para
comparar distribuciones entre clases, y se calculó la distancia de Cohen~$d$
para cuantificar la magnitud del efecto. El AUC univariante ---área bajo la
curva ROC usando cada variable como clasificador individual--- identificó las
20 features más discriminativas.

\begin{figure}[h]
\centering
%[FIGURA 1]
% Exportar desde Notebook 1_EDA.ipynb: gráfico de barras horizontales con el
% AUC univariante de las top 20 features, ordenado de mayor a menor.
% Ruta sugerida: figures/eda_auc_univariante.png
\includegraphics[width=\columnwidth]{figures/eda_auc_univariante.png}
\caption{AUC univariante de las 20 features con mayor poder discriminativo,
calculado sobre el dataset original de 50K sesiones. Las barras rojas
corresponden a features con AUC $\geq$ 0.80.}
\label{fig:eda_auc}
\end{figure}

Las cinco features más discriminativas resultaron ser
\texttt{phone\_call\_active}, \texttt{segmented\_typing\_ratio},
\texttt{hesitation\_count}, \texttt{data\_familiarity\_score} e
\texttt{input\_correction\_count}. Un Random Forest de validación rápida
obtuvo AUC$\approx$0.88 con validación cruzada de cinco pliegues,
confirmando la separabilidad del problema.

\begin{figure}[h]
\centering
%[FIGURA 2]
% Exportar desde Notebook 1_EDA.ipynb: panel de distribuciones bivariadas
% (legítimo vs. vishing) para las 6 features principales.
% Ruta sugerida: figures/eda_distribucion_features_clave.png
\includegraphics[width=\columnwidth]{figures/eda_distribucion_features_clave.png}
\caption{Distribución de las seis features más discriminativas para sesiones
legítimas (azul) y de vishing (rojo). Se aprecian diferencias marcadas en
tipeo segmentado, hesitación y monto de transacción.}
\label{fig:eda_dist}
\end{figure}

\subsection{Augmentación correlacionada (Fase 2)}

\subsubsection{Motivación}

SMOTE genera muestras sintéticas por interpolación lineal entre vecinos en el
espacio de features, lo que preserva distribuciones marginales pero no las
dependencias conjuntas entre variables~\cite{chawla2002smote}. En datos
comportamentales donde, por ejemplo, \texttt{call\_overlap\_duration\_s}
debe ser cero cuando \texttt{phone\_call\_active}$=0$, o donde
\texttt{hesitation\_count} y \texttt{avg\_hesitation\_duration\_s} deben
estar correlacionadas positivamente, esta limitación es especialmente dañina.

\subsubsection{CorrelatedAugmenter}

El \augcls{} implementa el proceso descrito en el Algoritmo~\ref{alg:aug}.
La idea central es generar ruido gaussiano multivariado con la misma
estructura de correlación que los datos originales, transformarlo a la escala
correcta mediante emparejamiento cuantil, y luego aplicar restricciones de
dominio para garantizar coherencia semántica.

\begin{algorithm}[h]
\caption{Generación de muestras con \augcls{}}
\label{alg:aug}
\begin{algorithmic}[1]
\REQUIRE Dataset de referencia $X$ por clase, número de muestras $n$,
         escala de ruido $\sigma = 0.03$
\STATE Separar $X$ en $X_{\text{legít.}}$ y $X_{\text{vishing}}$
\FOR{cada clase $c \in \{\text{legít.}, \text{vishing}\}$}
    \STATE Estimar la matriz de correlación $R_c$ sobre variables continuas
    \STATE Calcular la factorización de Cholesky $L_c = \text{chol}(R_c)$
    \STATE Estimar percentiles empíricos $\hat{F}_{c,j}$ para cada variable~$j$
    \STATE Generar ruido $Z \sim \mathcal{N}(0, I)$ de tamaño $n \times d$
    \STATE Obtener ruido correlacionado: $\tilde{Z} = Z \cdot L_c^\top$
    \STATE Transformar a $[0,1]$: $U = \Phi(\tilde{Z})$, donde $\Phi$ es la CDF normal
    \STATE Aplicar cuantil inverso: $\hat{X}_j = \hat{F}_{c,j}^{-1}(U_j)$
    \STATE Añadir perturbación: $\hat{X}_j \mathrel{+}= \mathcal{N}(0,\,\sigma \cdot \hat{\sigma}_j)$
    \STATE Aplicar restricciones de dominio (rangos, integridad semántica)
\ENDFOR
\RETURN Dataset sintético balanceado
\end{algorithmic}
\end{algorithm}

La validación de calidad del dataset de 1M sesiones mediante el test de
Kolmogorov-Smirnov comparado con el dataset original muestra una diferencia
media de~0.077 sobre todas las variables, lo que indica excelente preservación
de las distribuciones marginales. Solo dos variables presentan deriva leve:
\texttt{transaction\_amount\_cop} (KS$=0.24$) y \texttt{total\_dead\_time\_s}
(KS$=0.47$).

\begin{figure}[h]
\centering
%[FIGURA 3]
% Exportar desde Notebook 5_EDA_augmented_data.ipynb: gráfico de barras del
% estadístico KS por feature (original vs. aumentado), con línea de referencia
% en KS=0.10 (umbral de calidad aceptable).
% Ruta sugerida: figures/augmentacion_calidad_ks.png
\includegraphics[width=\columnwidth]{figures/augmentacion_calidad_ks.png}
\caption{Estadístico de Kolmogorov-Smirnov por variable: dataset original
(50K) vs.\ dataset aumentado (1M). Valores menores indican mejor preservación
de las distribuciones marginales. Media global: KS$=0.077$.}
\label{fig:ks}
\end{figure}

\subsection{Estrategias de balanceo de clases}

En ambas fases de modelado (Notebook~3 sobre 50K y Notebook~7 sobre 1M) se
aplican cuatro técnicas de re-muestreo en tres ratios de vishing objetivo
(10\,\%, 20\,\% y 25\,\%), generando 12~variantes de entrenamiento por tipo
de dato:

\begin{itemize}
    \item \textbf{Random Oversampling}: duplicación aleatoria de la clase
    minoritaria.
    \item \textbf{SMOTE}~\cite{chawla2002smote}: interpolación k-NN en el
    espacio de features.
    \item \textbf{Borderline-SMOTE}~\cite{han2005borderline}: SMOTE restringido
    a ejemplos en la frontera de decisión.
    \item \textbf{SMOTE + Undersampling}: reducción del 10\,\% de la clase
    mayoritaria combinada con SMOTE.
\end{itemize}

\subsection{Selección de métricas de evaluación}

En presencia de desbalance extremo (~1.5\,\% vishing en el dataset de 1M), el
\rocauc{} puede ser engañosamente optimista porque considera por igual los
verdaderos negativos y los verdaderos positivos. Con una tasa de vishing del
1.5\,\%, un clasificador trivial que etiquete todo como legítimo obtiene
ROC-AUC$\approx$0.5, pero en producción \emph{no detecta ningún fraude}.

El \prauc{} (área bajo la curva Precisión-Recall) es más informativo en este
contexto porque cuantifica directamente el trade-off entre detectar fraude
(Recall) y la calidad de las alertas generadas (Precisión)~\cite{davis2006pr}.
El umbral de clasificación se optimiza maximizando $F_1$ sobre la curva PR,
en lugar de usar el valor fijo de 0.5.

\subsection{Modelado multialgoritmo (Fase 2)}

Sobre los 12~datasets balanceados del dataset de 1M se entrenaron cuatro
algoritmos, evaluados sobre un holdout de 200K sesiones con distribución
real (~1.5\,\% vishing):

\begin{itemize}
    \item \textbf{Regresión Logística}: \texttt{max\_iter=1000}, como
    clasificador interpretable de referencia.
    \item \textbf{Random Forest}~\cite{breiman2001rf}: \texttt{n\_estimators=150},
    \texttt{max\_depth=10}.
    \item \textbf{XGBoost}~\cite{chen2016xgboost}: \texttt{max\_depth=6},
    \texttt{learning\_rate=0.1}, con aceleración GPU.
    \item \textbf{MLP}: dos capas ocultas (64-32 neuronas), 300~iteraciones.
\end{itemize}

\begin{figure}[h]
\centering
%[FIGURA 4]
% Exportar desde Notebook 7_Modeling_Vishing_AD_AWS_exec v2.ipynb: panel de
% barras comparando PR-AUC, Recall y F1 de los 4 algoritmos sobre los 12
% datasets balanceados (promediar por algoritmo o mostrar el mejor por algoritmo).
% Ruta sugerida: figures/nb7_comparacion_algoritmos.png
\includegraphics[width=\columnwidth]{figures/nb7_comparacion_algoritmos.png}
\caption{Comparación de PR-AUC por algoritmo y técnica de balanceo (Fase~2,
datos aumentados, holdout de 200K sesiones). XGBoost domina
consistentemente sobre todas las configuraciones.}
\label{fig:nb7_comp}
\end{figure}

XGBoost emergió como el algoritmo dominante en todas las configuraciones de
balanceo, motivando la exploración en profundidad de su espacio de
hiperparámetros en la Fase~3.

\subsection{Experimentación dirigida con XGBoost (Fase 3)}

\subsubsection{Variantes hiperparamétricas}

Se definieron siete variantes de XGBoost (Tabla~\ref{tab:xgb_variantes}) que
cubren diferentes ejes del espacio de hiperparámetros:

\begin{table}[h]
\caption{Variantes de XGBoost evaluadas en la Fase~3.}
\label{tab:xgb_variantes}
\centering
\begin{tabularx}{\columnwidth}{lX}
\toprule
\textbf{Variante} & \textbf{Descripción} \\
\midrule
\texttt{xgb\_base}        & Configuración idéntica a la Fase~2 (línea base) \\
\texttt{xgb\_deep}        & Árboles profundos (\texttt{max\_depth=9}), 300~estimadores, \texttt{lr=0.05} \\
\texttt{xgb\_shallow}     & Árboles poco profundos (\texttt{max\_depth=3}), 500~estimadores (boosting clásico) \\
\texttt{xgb\_regularized} & Regularización fuerte (L1=1.0, L2=5.0, \texttt{min\_child=10}, \texttt{gamma=0.3}) \\
\texttt{xgb\_balanced}    & \texttt{scale\_pos\_weight} calculado dinámicamente por dataset \\
\texttt{xgb\_conservative}& Subsampling agresivo (\texttt{subsample=0.7}, \texttt{colsample=0.7}, \texttt{gamma=0.5}) \\
\texttt{xgb\_slow\_learner}& Aprendizaje muy lento (\texttt{lr=0.01}, 500~estimadores) \\
\bottomrule
\end{tabularx}
\end{table}

La variante \texttt{xgb\_balanced} es la única cuyo hiperparámetro
\texttt{scale\_pos\_weight} varía por dataset: se calcula como la razón
$n_{\text{neg}}/n_{\text{pos}}$ del conjunto de entrenamiento, siguiendo la
recomendación del equipo de XGBoost para clases desbalanceadas~\cite{chen2016xgboost}.

\subsubsection{Cobertura del experimento}

Cada variante se entrenó sobre 13~datasets de datos aumentados (12~balanceados
más el raw de 1M) y 13~datasets de datos originales (12~balanceados más el raw
de 50K), con holdouts separados:

\[
7\,\text{variantes} \times 13\,\text{datasets} \times 2\,\text{tipos} = 182\,\text{modelos}
\]

Los modelos sobre datos originales se evaluaron sobre un holdout de 10K
sesiones (5\,\% vishing); los de datos aumentados, sobre 200K (1.5\,\% vishing).

% =============================================================================
\section{\wrapcls{}: Diseño para Producción}
\label{sec:wrapper}

Un requisito no funcional crítico del pipeline es la reproducibilidad de la
inferencia. En modelos de fraude, errores sutiles como un orden de features
diferente entre entrenamiento e inferencia, un scaler no guardado, o un umbral
fijo de 0.5 pueden invalidar toda la experiencia de entrenamiento.

\wrapcls{} es un patrón de empaquetado que encapsula en un único objeto
serializable con \texttt{joblib} los siguientes artefactos:

\begin{itemize}
    \item \textbf{Modelo entrenado}: instancia de \texttt{XGBClassifier}.
    \item \textbf{Lista de features}: orden exacto de las 44~variables.
    \item \textbf{Umbral óptimo}: calculado por maximización de $F_1$ en la
    curva PR sobre el holdout (no fijo en 0.5).
    \item \textbf{Metadatos}: nombre de variante, técnica de balanceo, ratio
    y tipo de dato.
\end{itemize}

La API expone tres métodos para satisfacer distintos casos de uso en
producción:

\begin{description}
    \item[\texttt{predict(json)}] Devuelve la etiqueta binaria (0/1).
    \item[\texttt{predict\_proba\_raw(json)}] Devuelve probabilidades
    \texttt{\{legitimate, vishing\}} para sistemas de ranking o gestión de colas.
    \item[\texttt{predict\_full(json)}] Devuelve etiqueta, probabilidades y
    umbral usado en un único diccionario.
\end{description}

El método acepta como entrada un \texttt{dict} Python, una cadena JSON o una
lista de dicts para procesamiento en batch. Valida features faltantes
explícitamente y lanza \texttt{ValueError} en lugar de fallar silenciosamente.

Los 182~modelos de la Fase~3 se almacenan en la ruta:
\begin{center}
\small
\texttt{s3://poc-fraude-vishing/proyecto/modelos\_xgb/\{tipo\}/\{variante\}/\{tecnica\}/\{ratio\}.pkl}
\end{center}

% =============================================================================
\section{Resultados}
\label{sec:resultados}

\subsection{Fase 1: Baseline sobre datos originales}

El modelado sobre el dataset de 50K sesiones con holdout de 10K (ratio real
5\,\% vishing) mostró que XGBoost con SMOTE+Undersampling al 10\,\% de ratio
de vishing objetivo obtuvo el mejor resultado (\prauc{}$\approx 0.40$--$0.50$).
La regresión logística, siendo el clasificador más simple, tuvo el peor
desempeño. El MLP se ubicó entre Logística y Random Forest, pero con mayor
varianza entre configuraciones de balanceo.

Estos resultados de línea base, aunque modestos, confirman la separabilidad
estadística del problema y orientan la necesidad de datos adicionales para
alcanzar métricas aptas para producción.

\subsection{Comparación de algoritmos en datos aumentados (Fase 2)}

Con el dataset de 1M sesiones y un holdout de 200K con desbalance real
(1.5\,\%), XGBoost supera consistentemente a los demás algoritmos en
\prauc{}, con ventaja más pronunciada en configuraciones de balanceo
conservadoras (10\,\% ratio de vishing objetivo). La Figura~\ref{fig:nb7_comp}
ilustra esta comparación.

\subsection{Experimentación XGBoost: variantes vs. técnicas de balanceo}

\begin{figure*}[t]
\centering
%[FIGURA 5]
% Exportar desde Notebook 9_XGBoost_training_exec.ipynb: heatmap doble
% (datos aumentados | datos originales) de PR-AUC promediado sobre ratios,
% con filas=técnica de balanceo y columnas=variante XGBoost.
% Ruta sugerida: figures/nb9_heatmap_pr_auc.png
\includegraphics[width=\textwidth]{figures/nb9_heatmap_pr_auc.png}
\caption{Heatmap de PR-AUC por variante de XGBoost y técnica de balanceo,
promediado sobre los tres ratios de vishing objetivo (10\,\%, 20\,\%, 25\,\%).
Izquierda: datos aumentados (1M). Derecha: datos originales (50K).}
\label{fig:heatmap}
\end{figure*}

La Figura~\ref{fig:heatmap} permite identificar los siguientes patrones:

\begin{enumerate}
    \item \texttt{xgb\_deep} y \texttt{xgb\_slow\_learner} dominan
    consistentemente sobre ambos tipos de datos.
    \item Random Oversampling al 25\,\% de ratio objetivo tiende a ser la
    mejor configuración de balanceo para datos originales.
    \item Los datos originales (50K) con balanceo adecuado superan a los datos
    aumentados (1M) en métricas de holdout, posiblemente porque el holdout de
    datos originales tiene un desbalance real del 5\,\% (más fácil) frente al
    1.5\,\% del holdout aumentado.
\end{enumerate}

\subsection{Mejor modelo consolidado}

La Tabla~\ref{tab:resultados} resume los mejores resultados por tipo de dato.
El mejor modelo global es \texttt{xgb\_deep} sobre datos originales con Random
Oversampling al 25\,\%.

\begin{table}[h]
\caption{Mejores resultados consolidados por tipo de dato (Fase~3, holdouts separados).}
\label{tab:resultados}
\centering
\begin{tabular}{lcccc}
\toprule
\textbf{Tipo dato} & \textbf{PR-AUC} & \textbf{Recall} & \textbf{Precisión} & \textbf{$F_1$} \\
\midrule
Original (50K)   & \textbf{0.9484} & 0.878 & 0.928 & 0.902 \\
Aumentado (1M)   & 0.8989          & 0.855 & 0.806 & 0.829 \\
\bottomrule
\multicolumn{5}{l}{\footnotesize Variante: \texttt{xgb\_deep} | Técnica: Random Oversampling 25\,\% | GPU: CUDA}
\end{tabular}
\end{table}

\begin{figure}[h]
\centering
%[FIGURA 6]
% Exportar desde Notebook 9_XGBoost_training_exec.ipynb: curva PR del mejor
% modelo (xgb_deep / original / random_oversampling / 25%), con el punto
% óptimo F1 marcado y el área bajo la curva sombreada.
% Ruta sugerida: figures/nb9_curva_precision_recall.png
\includegraphics[width=\columnwidth]{figures/nb9_curva_precision_recall.png}
\caption{Curva Precisión-Recall del mejor modelo (\texttt{xgb\_deep},
datos originales, Random Oversampling~25\,\%). PR-AUC$=0.9484$; punto
óptimo $F_1=0.902$ en umbral$\approx$0.XX (completar con el valor del notebook).}
\label{fig:curva_pr}
\end{figure}

\begin{figure}[h]
\centering
%[FIGURA 7]
% Exportar desde Notebook 9_XGBoost_training_exec.ipynb: gráfico de barras
% horizontales de feature importance (gain) del mejor modelo xgb_deep.
% Top 20 features, coloreadas por cuartil de importancia.
% Ruta sugerida: figures/nb9_feature_importance_xgb_deep.png
\includegraphics[width=\columnwidth]{figures/nb9_feature_importance_xgb_deep.png}
\caption{Importancia de features (gain de XGBoost) del mejor modelo
\texttt{xgb\_deep}. Las barras rojas corresponden al top~20\,\% de
importancia. \texttt{phone\_call\_active} y \texttt{hesitation\_count}
encabezan el ranking.}
\label{fig:feat_imp}
\end{figure}

% =============================================================================
\section{Simulación de Inferencia en Tiempo Real}
\label{sec:simulacion}

\subsection{Configuración}

Para validar el pipeline de punta a punta se generó un dataset independiente
de 100\,000~sesiones 100\,\% sintéticas mediante \augcls{} entrenado sobre el
dataset original de 50K como referencia estadística (ninguna observación del
entrenamiento se incluye en el dataset de simulación). La distribución objetivo
es 1.5\,\% vishing (98\,500 legítimas, 1\,500 de vishing), consistente con la
prevalencia real del dataset aumentado.

El modelo seleccionado para la simulación es \texttt{xgb\_deep} entrenado sobre
datos aumentados con Random Oversampling al 25\,\%
(\texttt{threshold}$=0.9365$). Se eligió la versión de datos aumentados para
la simulación porque el holdout de 200K refleja mejor la distribución de
producción (1.5\,\% vishing), a diferencia del holdout de datos originales
(5\,\%).

Cada observación se envía \emph{individualmente} como \texttt{dict} al
\wrapcls{} y la latencia se mide con \texttt{time.perf\_counter()},
replicando el flujo de inferencia sesión a sesión en producción.

\subsection{Resultados de detección}

\begin{table}[h]
\caption{Métricas de detección en la simulación de inferencia (100K sesiones).}
\label{tab:simulacion}
\centering
\begin{tabular}{lc}
\toprule
\textbf{Métrica} & \textbf{Valor} \\
\midrule
Sesiones vishing reales (TP + FN) & 1\,500 \\
Sesiones legítimas (TN + FP)     & 98\,500 \\
\midrule
Verdaderos positivos (TP)    & 1\,133 \\
Falsos negativos (FN)        & 367    \\
Falsos positivos (FP)        & 337    \\
Verdaderos negativos (TN)    & 98\,163 \\
\midrule
Recall                       & 0.7553  \\
Precisión                    & 0.7707  \\
$F_1$                        & 0.7630  \\
Tasa de alerta               & 1.47\,\% \\
\midrule
\textbf{Latencia media}      & \textbf{1.5978~ms} \\
Latencia mediana (p50)       & 1.5273~ms \\
Latencia p95                 & 1.7939~ms \\
Latencia p99                 & 2.2447~ms \\
Latencia máxima              & 49.45~ms  \\
\textbf{Rendimiento}         & \textbf{539~obs/s} \\
\bottomrule
\end{tabular}
\end{table}

\begin{figure}[h]
\centering
%[FIGURA 8]
% Exportar desde Notebook 10_Realtime_Simulation.ipynb (celda 4.2):
% Panel de 3 gráficas:
%   - Izquierda: distribución de scores separada por clase real (legítimo/vishing)
%     con línea vertical en el umbral (0.9365)
%   - Centro: matriz de confusión (conteos)
%   - Derecha: matriz de confusión normalizada por fila
% Ruta sugerida: figures/realtime_confusion_scores.png
\includegraphics[width=\columnwidth]{figures/realtime_confusion_scores.png}
\caption{Distribución de scores de vishing por clase real (izquierda) y
matrices de confusión absolutas y normalizadas (centro y derecha) para la
simulación sobre 100K sesiones. Umbral óptimo: 0.9365.}
\label{fig:sim_confusion}
\end{figure}

La diferencia en $F_1$ entre el holdout de entrenamiento ($F_1=0.829$ para
datos aumentados) y la simulación ($F_1=0.763$) refleja principalmente que
el dataset de simulación fue generado de forma completamente independiente
---sin solapamiento de filas con el entrenamiento--- y con un \augcls{}
reentrenado desde cero, lo que elimina cualquier solapamiento implícito de
distribuciones.

\subsection{Análisis de latencia}

\begin{figure}[h]
\centering
%[FIGURA 9]
% Exportar desde Notebook 10_Realtime_Simulation.ipynb (celda 4.1):
% Panel de 3 gráficas: histograma de latencias, boxplot por clase predicha,
% y curva de percentiles acumulados.
% Ruta sugerida: figures/realtime_latencia_distribucion.png
\includegraphics[width=\columnwidth]{figures/realtime_latencia_distribucion.png}
\caption{Distribución de latencias de inferencia sobre 100K observaciones.
Izquierda: histograma con marcas en p50, p95 y p99. Centro: boxplot por
clase predicha. Derecha: curva de percentiles acumulados.}
\label{fig:latencia}
\end{figure}

La latencia media de 1.60~ms y el percentil~99 de 2.24~ms son consistentes
con un despliegue en producción dentro del ciclo de respuesta de la
aplicación bancaria, donde las sesiones típicas duran decenas de segundos.
El pico de 49.45~ms corresponde al primer ciclo de inferencia (JIT
compilation de XGBoost sobre CUDA) y no es representativo del
comportamiento estacionario. El rendimiento de 539~obs/s cubre ampliamente
la carga transaccional de referencia de la aplicación.

El análisis del timeline de riesgo (score de vishing promediado en ventana
móvil de 100~observaciones) muestra que las sesiones de vishing presentan
clusters temporales distinguibles sobre el fondo de sesiones legítimas,
consistente con la hipótesis de que el comportamiento inducido por el atacante
es cualitativamente diferente al comportamiento orgánico.

% =============================================================================
\section{Discusión}
\label{sec:discusion}

\subsection{Features más importantes e implicaciones}

Las cinco features con mayor importancia en el modelo final son consistentes
con las hipótesis comportamentales del vishing:

\begin{enumerate}
    \item \textbf{\texttt{phone\_call\_active}}: la presencia confirmada de
    una llamada activa es el indicador más directo. La aplicación bancaria
    accede a este estado a través de las APIs del sistema operativo.

    \item \textbf{\texttt{hesitation\_count}} y
    \textbf{\texttt{avg\_hesitation\_duration\_s}}: el usuario en vishing
    duda consistentemente antes de cada acción porque espera instrucciones
    del atacante.

    \item \textbf{\texttt{segmented\_typing\_ratio}}: el tipeo segmentado
    (escribir campo a campo, dictando lo que escucha) es un patrón marcado
    frente al tipeo fluido del usuario en sesión legítima.

    \item \textbf{\texttt{input\_correction\_count}}: el usuario comete más
    errores al transcribir datos dictados (número de cuenta, monto) que al
    ingresar datos memorizados.
\end{enumerate}

Estos hallazgos sugieren que un modelo más simple, incluso una regresión
logística con las 10 features más importantes, podría ser suficiente para
una primera versión en producción con mayor interpretabilidad.

\subsection{Impacto del threshold dinámico}

El umbral óptimo del modelo usado en simulación fue~0.9365, muy superior al
0.5 convencional. Esto refleja que el modelo asigna probabilidades de vishing
bajas a la mayoría de sesiones (clase dominante), y solo supera el umbral para
casos con señales comportamentales muy marcadas. Un umbral fijo en 0.5
generaría una tasa de falsas alertas inaceptablemente alta en producción.

\subsection{Limitaciones del estudio}

\begin{enumerate}
    \item \textbf{Dataset sintético}: aunque el \augcls{} preserva la
    estructura estadística, las distribuciones de los datos reales de BioCatch
    son desconocidas. La validación sobre datos reales etiquetados es el paso
    crítico pendiente antes de cualquier despliegue.

    \item \textbf{Deriva temporal}: el modelo fue entrenado y evaluado sobre
    el mismo período simulado. En producción, los patrones de vishing evolucionan
    y el modelo requeriría re-entrenamiento periódico~\cite{lu2020concept}.

    \item \textbf{Variabilidad del holdout}: el gap entre el $F_1=0.902$
    del holdout de entrenamiento y el $F_1=0.763$ de la simulación
    independiente indica que parte del rendimiento se debe a la homogeneidad
    estadística del dataset sintético. Con datos reales, esta diferencia
    podría ser mayor.

    \item \textbf{Sin exploración de SHAP}: el análisis de importancia actual
    usa el gain de XGBoost, que no es aditivo. Un análisis de valores
    SHAP~\cite{lundberg2017shap} permitiría interpretaciones más robustas y
    por instancia.

    \item \textbf{Latencia medida en SageMaker}: las cifras de latencia
    corresponden a un entorno de cómputo acelerado en nube. En infraestructura
    de producción local o en contenedores sin GPU, la latencia podría ser
    superior.
\end{enumerate}

\subsection{Trabajo futuro}

\begin{itemize}
    \item \textbf{Validación con datos reales}: coordinar con BioCatch o
    equipos de fraude para obtener un dataset real etiquetado y validar el
    pipeline.

    \item \textbf{Feature adicional}: \texttt{background\_switch\_count}
    (número de veces que el usuario sale de la app para consultar información),
    identificada como señal proximal de vishing cuando \texttt{phone\_call\_active}
    no está disponible.

    \item \textbf{Modelos online}: explorar actualización incremental del modelo
    (\textit{online learning}) para adaptarse a la deriva temporal sin reentrenamiento
    completo~\cite{lu2020concept}.

    \item \textbf{Interpretabilidad con SHAP}: calcular valores SHAP globales
    y por instancia para mejorar la explicabilidad de las alertas generadas
    ante equipos de cumplimiento y auditoría~\cite{lundberg2017shap}.

    \item \textbf{Integración con sistema de alertas}: definir el flujo de
    gestión de alertas (cola de revisión humana, interrupción automática de
    sesión, notificación al cliente) y medir el impacto operacional.

    \item \textbf{Pruebas adversariales}: evaluar la robustez del modelo ante
    atacantes que conozcan las señales capturadas e intenten imitar el
    comportamiento legítimo.
\end{itemize}

% =============================================================================
\section{Conclusiones}
\label{sec:conclusiones}

Este trabajo demuestra la viabilidad técnica de un sistema de detección de
vishing basado en biometría comportamental, entrenado íntegramente sobre datos
sintéticos. Los resultados principales son:

\begin{enumerate}
    \item El problema es estadísticamente separable: incluso con datos
    sintéticos y un modelo simple, el AUC univariante de las features más
    discriminativas supera~0.80.

    \item El \augcls{} permite escalar de 50K a 1M sesiones preservando la
    estructura de correlación multivariada, con una fidelidad medida por
    KS$=0.077$ de media.

    \item \texttt{xgb\_deep} con Random Oversampling al 25\,\% es la
    configuración dominante, alcanzando PR-AUC$=0.9484$, Recall$=0.878$ y
    $F_1=0.902$ en holdout de datos originales.

    \item La simulación de inferencia en tiempo real valida la viabilidad
    operacional: Recall$=0.755$, $F_1=0.763$, latencia media~1.60~ms y
    p99~$<$~2.25~ms con un rendimiento de~539~obs/s.

    \item El patrón \wrapcls{} resuelve el problema de reproducibilidad de
    inferencia al encapsular todos los artefactos necesarios en un único
    objeto serializable.
\end{enumerate}

El paso crítico siguiente es la validación sobre datos reales etiquetados.
Los supuestos comportamentales del dataset sintético están alineados con
la literatura de BioCatch, pero la calibración final del modelo debe hacerse
sobre el comportamiento real de los usuarios de la aplicación bancaria.

% =============================================================================
\section*{Agradecimientos}

Los autores agradecen al equipo de [Área/División] por el soporte en
infraestructura AWS y a BioCatch por la documentación pública sobre
biometría comportamental que inspiró el diseño del dataset.

% =============================================================================
\bibliographystyle{IEEEtran}
\begin{thebibliography}{99}

\bibitem{mitnick2002art}
K.~D. Mitnick y W.~L. Simon,
\textit{The Art of Deception: Controlling the Human Element of Security}.
Indianapolis: Wiley Publishing, 2002.

\bibitem{jagatic2007social}
T.~N. Jagatic, N.~A. Johnson, M.~Jakobsson, y F.~Menczer,
``Social phishing,''
\textit{Communications of the ACM}, vol.~50, no.~10, pp.~94--100, oct. 2007.

\bibitem{threatfabric2023}
ThreatFabric,
``Mobile Banking Fraud Report 2023,''
ThreatFabric B.V., Reporte técnico, 2023.
[En línea]. Disponible: \url{https://www.threatfabric.com/blogs/mobile-banking-fraud}

\bibitem{biocatch2022}
BioCatch,
``The Science of Behavioral Biometrics: Detecting Social Engineering and
Account Takeover Fraud,''
BioCatch Ltd., Libro blanco, 2022.
[En línea]. Disponible: \url{https://www.biocatch.com}

\bibitem{shen2007application}
A.~Shen, R.~Tong, y Y.~Deng,
``Application of classification models on credit card fraud detection,''
en \textit{Proc. Int. Conf. Service Systems and Service Management (ICSSSM)},
Chengdu, China, 2007, pp.~1--4.

\bibitem{sundararajan2020behavioral}
K.~Sundararajan y T.~Han,
``Behavioral biometrics: A survey on human recognition through activity
patterns on smartphones,''
\textit{IEEE Access}, vol.~8, pp.~80837--80867, 2020.

\bibitem{chen2016xgboost}
T.~Chen y C.~Guestrin,
``XGBoost: A scalable tree boosting system,''
en \textit{Proc. 22nd ACM SIGKDD Int. Conf. Knowledge Discovery and Data Mining
(KDD)}, San Francisco, CA, 2016, pp.~785--794.

\bibitem{phua2010comprehensive}
C.~Phua, V.~Lee, K.~Smith, y R.~Gayler,
``A comprehensive survey of data mining-based fraud detection research,''
\textit{arXiv preprint arXiv:1009.6119}, 2010.

\bibitem{bhattacharyya2011data}
S.~Bhattacharyya, S.~Jha, K.~Tharakunnel, y J.~C. Westland,
``Data mining for credit card fraud: A comparative study,''
\textit{Decision Support Systems}, vol.~50, no.~3, pp.~602--613, feb. 2011.

\bibitem{chawla2002smote}
N.~V. Chawla, K.~W. Bowyer, L.~O. Hall, y W.~P. Kegelmeyer,
``SMOTE: Synthetic minority over-sampling technique,''
\textit{Journal of Artificial Intelligence Research}, vol.~16, pp.~321--357,
jun. 2002.

\bibitem{han2005borderline}
H.~Han, W.~Y. Wang, y B.~H. Mao,
``Borderline-SMOTE: A new over-sampling method in imbalanced data sets
learning,''
en \textit{Proc. Int. Conf. Intelligent Computing (ICIC)}, Hefei, China,
2005, pp.~878--887.

\bibitem{he2009learning}
H.~He y E.~A. Garcia,
``Learning from imbalanced data,''
\textit{IEEE Transactions on Knowledge and Data Engineering}, vol.~21, no.~9,
pp.~1263--1284, sep. 2009.

\bibitem{krawczyk2016learning}
B.~Krawczyk,
``Learning from imbalanced data: Open challenges and future directions,''
\textit{Progress in Artificial Intelligence}, vol.~5, no.~4, pp.~221--232,
oct. 2016.

\bibitem{davis2006pr}
J.~Davis y M.~Goadrich,
``The relationship between precision-recall and ROC curves,''
en \textit{Proc. 23rd Int. Conf. Machine Learning (ICML)}, Pittsburgh, PA,
2006, pp.~233--240.

\bibitem{breiman2001rf}
L.~Breiman,
``Random forests,''
\textit{Machine Learning}, vol.~45, no.~1, pp.~5--32, oct. 2001.

\bibitem{patki2016synthetic}
N.~Patki, R.~Wedge, y K.~Veeramachaneni,
``The synthetic data vault,''
en \textit{Proc. IEEE Int. Conf. Data Science and Advanced Analytics (DSAA)},
Montreal, Canada, 2016, pp.~399--410.

\bibitem{nikolentzos2021synthetic}
G.~Nikolentzos, G.~Bouritsas, y M.~Vazirgiannis,
``Synthetic data generation for fraud detection using graph neural networks,''
\textit{arXiv preprint arXiv:2106.09399}, 2021.

\bibitem{fridman2017active}
L.~Fridman, S.~Weber, R.~Greenstadt, y M.~Kam,
``Active authentication on mobile devices via stylometry, application usage,
web browsing, and GPS location,''
\textit{IEEE Systems Journal}, vol.~11, no.~2, pp.~513--521, jun. 2017.

\bibitem{tetri2013dissecting}
P.~Tetri y J.~Vuorinen,
``Dissecting social engineering,''
\textit{Behaviour \& Information Technology}, vol.~32, no.~10, pp.~1014--1023,
2013.

\bibitem{apwg2023}
Anti-Phishing Working Group (APWG),
``Phishing Activity Trends Report, Q4 2023,''
APWG, Reporte, 2023.
[En línea]. Disponible: \url{https://apwg.org/trendsreports/}

\bibitem{lu2020concept}
J.~Lu, A.~Liu, F.~Dong, F.~Gu, J.~Gama, y G.~Zhang,
``Learning under concept drift: A review,''
\textit{IEEE Transactions on Knowledge and Data Engineering}, vol.~31, no.~12,
pp.~2346--2363, dic. 2019.

\bibitem{lundberg2017shap}
S.~M. Lundberg y S.-I. Lee,
``A unified approach to interpreting model predictions,''
en \textit{Advances in Neural Information Processing Systems (NeurIPS)},
vol.~30, Long Beach, CA, 2017.

\bibitem{lemaitre2017imbalanced}
G.~Lema\^{i}tre, F.~Nogueira, y C.~K. Aridas,
``Imbalanced-learn: A Python toolbox to tackle the curse of imbalanced
datasets in machine learning,''
\textit{Journal of Machine Learning Research}, vol.~18, no.~17, pp.~1--5,
2017.

\end{thebibliography}

\end{document}
